Under the stricter sibling-group split, it remains above G-Memory at 72B (31.9 vs.\ 14.4\%, $p=6.4\!\times\!10^{-25}$) and 30B (15.9 vs.\ 11.4\%, $p=3.2\!\times\!10^{-5}$), but not at 7B (3.7 vs.\ 3.3\%, $p=.665$). The 30B and 7B comparison baselines use the think condition, and no matched no-think arm is available for ours.

\begin{figure}[H]
    \centering
    \includegraphics[width=\linewidth]{figs/viki_main_results.pdf}
    \caption{VIKI-L2 success. OOD denotes single-family withholding, and CG denotes compositional generalization with or without image input. Numeric labels mark our method. Sibling-group-held-out OOD appears in Appendix~\ref{sec:viki_audit}.}
    \label{fig:viki_main}
\end{figure}

ToM and zero-shot both emit primitive plans, and ToM differs significantly from zero-shot in none of the 12 paired cells (all $p\geq.296$). In its highest-scoring cell, 72B OOD, ToM raises in-scene-object validity from 52.7\% to 64.4\%, yet success moves from 3.5\% to 2.7\% (Table~\ref{tab:tom_funnel}, Figure~\ref{fig:tom_failure}). Our method instead outputs skill calls, so its gap to the primitive-plan baselines measures the skill-call interface and the library together. The RQ2 arms hold the interface and the planner fixed and change only the library. The no-trace library yields zero successes in every cell, the library without admission averages 21.9\% at 72B, and the admitted library averages 74.6\% (Section~\ref{subsec:induction_validation}). The interface alone therefore produces no success, and within it the library accounts for the gain. With the admitted library placed in the prompt and primitive plans as output, success on the two CG splits is 0.0\% with images and 3.0\% without, not above G-Memory under the same output format (Appendix~\ref{sec:viki_audit}). On VIKI-L2 the gain therefore requires the planner that consumes the library, and we treat the planner as part of the method.

\subsection{Validating Trace-Based Skill Induction (RQ2)}
\label{subsec:induction_validation}

[... unchanged text between excerpts ...]

\paragraph{Limitations.} We study alternating rather than concurrent actions. Each VIKI-L2 library is built once, and the 30B full arm has one run and no matched no-think arm. The PARTNR libraries behind Tables~\ref{tab:main_results_task} and~\ref{tab:partnr_composition} are built without the execution gate, and the PARTNR gate validation uses privileged state. On PARTNR the partner condition of a cooperation skill is free text that only the planner reads, and observed state changes trigger replanning but are not written into the prompt. The VIKI-L2 CG split has one composition pattern, and no single setting combines decentralized cooperation with a well-powered compositional test. On VIKI-L2 the library gains over G-Memory only when the deterministic planner consumes it. Low 7B sibling-group success precludes transfer claims for smaller backbones.
% Acknowledgments are omitted in the anonymous submission.  Add an
% acknowledgments paragraph here for the camera-ready version.

[... unchanged text between excerpts ...]

\paragraph{CG split audit.}
A read-only audit compares the CG tasks with every episode pool used to build the library evaluated in Figure~\ref{fig:viki_main}, namely the seeds, holdouts, and tool calls of its 896 induction runs, the support probe, and the 3{,}598-episode pool from which the ordering rules are mined. All 1{,}242 training episodes that the proposing agent addresses through tools lie in the induction half, and none of the 896 transcripts contains a CG task identifier. Four of the 295 CG task identifiers also label training episodes, two of them in the induction half, because VIKI-L2 reuses identifiers across its training and test files. Each paired episode differs from its CG task in instruction, initial positions, goals, and reference plan, and none contains the held-out combination. The 297 CG rows hold 295 distinct tasks because two tasks occur twice. The unseen property is the combination of a two-robot cutting subtask with an independent delivery subtask. A weaker structure, an activation step followed by a transport step, does occur in training (1{,}033 rows), so we do not claim that this structure is unseen.

\paragraph{In-context library control.}
This arm places the admitted library in the prompt of Qwen2.5-VL-72B-Instruct and asks for a primitive plan. It keeps the G-Memory prompt, model, decoding, and scoring of Figure~\ref{fig:viki_main} and replaces only the memory block with one fixed rendering of the eight operators and the three ordering rules, identical for every row. The G-Memory and zero-shot responses are those of Figure~\ref{fig:viki_main}. On CG w/o Image the arm succeeds on 9 of 297 rows under both the JSON-tolerant parser and the official scorer, against 10 for G-Memory (9 vs.\ 10 rows solved by only one arm, $p=1.0$) and 1 for zero-shot (8 vs.\ 0, $p=.0078$). On CG w/ Image it succeeds on no row under either criterion, against 14 for G-Memory (0 vs.\ 14, $p=1.2\times10^{-4}$) and none for zero-shot. Our method solves 247 and 237 rows that this arm misses and loses none ($p<10^{-70}$ on each split). Dropping the second occurrence of the two repeated tasks changes neither direction nor significance. Rerunning the first 20 rows of CG w/o Image changes 19 of the 20 responses and their success count from 0 to 3, so single-run differences of a few rows are within run-to-run variation.

\paragraph{Sibling-group-held-out OOD.}
